% ═══════════════════════════════════════════════════════════════════════════
%  INSTÄLLNINGAR — Reformationsbibeln SRB 2026
%
%  Allt som styr bokens utseende ligger i denna fil. main.tex innehåller
%  bara logiken och ska normalt inte behöva ändras.
%
%  Efter ändring: kompilera om (xelatex main.tex, fyra gånger för hela bibeln).
%  Ångra allt:    python3 versioner.py hamta <tagg>
%  Alla val och sidkostnader finns beskrivna i PARAMETRAR.md
% ═══════════════════════════════════════════════════════════════════════════


% ── 1. TYPOGRAFI ──────────────────────────────────────────────────────────
% Sidkostnad för hela bibeln inom parentes.

\def\SRBgrad{10.5pt}          % teckenstorlek        (13pt ≈ −70 sidor, 15pt ≈ +70)
\def\SRBradavstand{14pt}     % radavstånd           (±22 sidor per 0,4pt)
\def\SRBfetning{1.26}        % FakeBold 0–1.8       (0 sidor — ändrar inte radbrytning)
\def\SRBsiffror{OldStyle}   % OldStyle = gemena siffror, Lining = versala
\def\SRBtolerance{3000}     % radbrytningens tolerans; högre = tätare ordmellanrum


% ── 2. ANFANG (stora begynnelsebokstaven) ─────────────────────────────────

\def\SRBanfangrader{2}      % antal rader hög      (3 ≈ +26 sidor, 1 ≈ −26)
\def\SRBanfangJ{0.18}       % storlekjustering för bokstaven J
\def\SRBanfangJlyft{0.15}
\def\SRBanfangOvriga{0.25}  % storlekjustering för övriga bokstäver
\def\SRBanfangOvrigalyft{0.01}
\def\SRBanfangKapitaler{1}  % 1 = första ordet i kapitäler, 0 = vanlig stil


% ── 3. TEXTENS STRUKTUR ───────────────────────────────────────────────────

\def\SRBkapitelminplats{1}  % rader som minst måste rymmas innan en kapitelrubrik
                            % (annars hänger anfangen ned i not- och sidfotszonen)
\def\SRBkapitelluft{0.0cm} % mellanrum mellan kapitel   (±15 sidor per 0,1cm)
\def\SRBbokluft{5mm}        % mellanrum mellan böcker på samma sida
\def\SRBnybokgrans{0mm}    % ny bok får ny sida om mindre än detta återstår
                            % sätt 300mm => alltid ny sida (≈ +25 sidor)
                            % sätt 0mm   => aldrig ny sida
\def\SRBankor{10000}        % 10000 = skydda mot änkor/horungar, 300 = tillåt
\def\SRBindrag{0em}       % styckeindrag


% ── 4. SIDLAYOUT — måtten sätts av format/A4.tex eller format/BOK.tex ─────────────────────────────────────────────────────
% Sidkostnad: ~11 sidor per 2mm över/under, ~6 sidor per 2mm inner/ytter.

% \def\SRBpappersbredd{210mm}   <- sätts av format/<FORMAT>.tex
% \def\SRBpappersbreddMM{210}  % samma mått utan enhet (för sidnummer och tumindex)   <- sätts av format/<FORMAT>.tex
% \def\SRBpappershojd{297mm}   <- sätts av format/<FORMAT>.tex
% \def\SRBpappershojdMM{297}   <- sätts av format/<FORMAT>.tex
% \def\SRBover{1mm}         % övermarginal   <- sätts av format/<FORMAT>.tex
% \def\SRBunder{15mm}        % undermarginal   <- sätts av format/<FORMAT>.tex
% \def\SRBinner{14mm}        % innermarginal (ryggsidan)   <- sätts av format/<FORMAT>.tex
% \def\SRBytter{10mm}        % yttermarginal   <- sätts av format/<FORMAT>.tex
% \def\SRBfotavstand{1.2cm}   % avstånd ned till sidfoten   <- sätts av format/<FORMAT>.tex


% ── 5. SIDHUVUD OCH SIDFOT ────────────────────────────────────────────────

\def\SRBsidnummerstorlek{11pt}
\def\SRBsidnummerTopp{0}       % gammalt sidhuvud, av (numret klipptes vid kanten)
\def\SRBsidnummerMarginal{}   % V17: sidnumret i yttermarginalen i stället
% \def\SRBsidnummerKant{6.5}       % mm från papperets ytterkant   <- sätts av format/<FORMAT>.tex
\def\SRBsidnummerMarginalstorlek{6pt} % sidnumrets storlek i marginalen
\def\SRBsidnummerHojd{3}       % mm från papperets överkant (ovanför textblocket)
\def\SRBsidnummerBotten{1}  % 1 = sidnummer centrerat nederst, 0 = av

% Vad som står i sidfotens ytterkant:
%   kapitel  = "Första Moseboken 12 – Första Moseboken 14"  (nuvarande)
%   position = "Kapitel 3 av 50"
%   perikop  = perikopens namn        (kräver SRBperikoprubriker=1)
%   inget    = tomt
\def\SRBsidfot{inget}
\def\SRBsidfotEndastUtanRubrik{0}  % 1 = visa bara på sidor utan kapitelrubrik


% ── 6. TUMINDEX ───────────────────────────────────────────────────────────

\def\SRBframsida{0}         % 1 = titelsidor och innehållsförteckning med,
                            % 0 = hoppa över dem (snabba provbyggen)
\def\SRBtumindex{0}         % 1 = på, 0 = av
\def\SRBtabbredd{9}         % mm flikbredd
\def\SRBtabminhojd{1.6}     % mm minsta flikhöjd (så korta böcker syns)
\def\SRBtablucka{3.2}       % mm lucka mellan de sex avdelningarna
\def\SRBtabover{18}         % mm ovanför första fliken
\def\SRBtabunder{16}        % mm nedanför sista fliken


% ── 7. FÄRG ───────────────────────────────────────────────────────────────
% Gråskala (0 = svart, 1 = vitt). Krävs för tryck — en tryckplåt.

\def\SRBtextton{0}          % brödtext
\def\SRBtillaggston{0.50}   % tilläggsord som ej finns i grundtexten
                            % mörkare (0.38) syns bättre på gräddpapper
\def\SRBrubrikstil{kursiv}  % kursiv eller sparrad (upprätt, breddad) - LÅST som standard
\def\SRBrubrikAndel{1.05}    % V41: mycket mindre, kompakt - nästan brödtextens storlek
\def\SRBrubrikfetning{0.4}   % V41: LITE bold - diskret tyngd, inte kraftig
\def\SRBrubrikton{0.1}     % rubriker och innehållsförteckningens sidnummer
\def\SRBbakgrund{1}         % sidbakgrund; 1 = vitt (papperet ger gräddtonen)


% ── 8. INNEHÅLL SOM FINNS I TEXTEN MEN KAN DÖLJAS ────────────────────────
% Allt nedan GENERERAS ALLTID ut i textfilerna. Här styrs bara om det VISAS.
% usx_config.json behöver inte längre ändras.

\def\SRBversnummer{1}          % versnummer i texten (0 sidor - \smash-tekniken gör dem gratis)
\def\SRBversnummerAndel{0.46}     % versnumrets storlek som andel av brödtexten
\def\SRBversnummerLyftAndel{0.45} % hur högt numret lyfts, andel av brödtexten
\def\SRBversnummerstorlek{7.5pt}  % (används ej längre - kvar för bakåtkompat.)
\def\SRBversnummerHoger{0.09em}   % liten luft EFTER numret, innan källans ~ (hårt mellanslag)
\def\SRBversnummerVanster{-0.05em} % liten luft FÖRE numret, utöver källans egna mellanslag
\def\SRBversnummerLyft{1.45ex}    % hur högt numret lyfts över raden

% Noternas utseende. De sätts som ett löpande stycke längst ned på sidan,
% direkt ovanför sidfoten, och tar därför normalt bara en rad.
\def\SRBnotstil{\fontsize{8.8pt}{10pt}\selectfont}
\def\SRBnotmarkestorlek{7pt} % hänvisningssiffrans storlek (samma för not och korsref)
\def\SRBnotmarketon{0.30}    % hänvisningssiffrans gråton (var 0.45 — nu fylligare)
\def\SRBnotmarkefet{1}       % 1 = fet siffra i notblocket    % hänvisningssiffrans gråton; högre = ljusare
\def\SRBnotavstand{7pt}      % luft mellan brödtext och noter
\def\SRBnotlinje{0.25\textwidth} % skiljelinjens längd
\def\SRBfotnoter{1}            % TR/UN-varianter i sidfoten    (≈ +100–150 sidor)
\def\SRBkorshanvisningar{0}    % parallellställen i sidfoten   (≈ +60–100 sidor)

% Perikopbrytpunkterna: samma ställe i texten, tre olika sätt att visa det.
\def\SRBperikoprubriker{0}     % rubrik i texten               (≈ +15–25 sidor)
\def\SRBstyckebrytning{1}      % nytt stycke, ingen rubrik     (≈ +65 sidor, valt)
\def\SRBstyckemarkor{0}        % grått ¶ utan styckebrytning   (0 sidor)
\def\SRBsidfotPerikopMark{0}   % spara perikopnamn för sidfoten

\def\SRBsuperGray{0.50}